\subsection{GROUP ALGEBRA OF A FREE ABELIAN GROUP}

Let P be a free Z-module of finite rank l. We denote by A{[}P{]} the group algebra of the additive group of P over A (Algebra, Chap. III, § 2, no. 6). For any \(p \in P\) , denote by \(e^{p}\) the corresponding element of A{[}P{]}. Then \((e^{p})_{p \in P}\) is a basis of the A-module A{[}P{]}, and, for any \(p, p' \in P\) , we have

\[
e ^ {p} e ^ {p ^ {\prime}} = e ^ {p + p ^ {\prime}}, \quad (e ^ {p}) ^ {- 1} = e ^ {- p}, \quad e ^ {0} = 1.
\]

Lemma 1. Assume that A is factorial (Commutative Algebra, Chap. VII, § 3, no. 1, Def. 1).

\begin{enumerate}
\def\labelenumi{(\roman{enumi})}
\item
  The ring A{[}P{]} is factorial.
\item
  If \(u, v\) are non-proportional elements of \(\mathbf{P}\) , the elements \(1 - e^u, 1 - e^v\) of \(\mathrm{A}[\mathbf{P}]\) are relatively prime.
\end{enumerate}

Let \((p_{1}, p_{2}, \ldots, p_{l})\) be a basis of P, and \(X_{1}, X_{2}, \ldots, X_{l}\) be indeterminates. The A-linear map from \(A[X_{1}, \ldots, X_{l}, X_{1}^{-1}, \ldots, X_{l}^{-1}]\) to A{[}P{]} that takes \(X_{1}^{n_{1}} X_{2}^{n_{2}} \ldots X_{l}^{n_{l}}\) (where \(n_{1}, n_{2}, \ldots, n_{l} \in Z\) ) to \(e^{n_{1} p_{1} + \cdots + n_{l} p_{l}}\) is an isomorphism of rings. Now \(A[X_{1}, \ldots, X_{l}]\) is a factorial ring (Commutative Algebra, Chap. VII, § 3, no. 5), and \(A[X_{1}, \ldots, X_{l}, X_{1}^{-1}, \ldots, X_{l}^{-1}]\) is the ring of fractions of \(A[X_{1}, \ldots, X_{l}]\) , hence is also factorial.

Let \(P'\) (resp. \(P''\) ) be the set of elements of \(P\) of which some multiple belongs to \(Zu + Zv\) (resp. \(Zu\) ). Then the groups \(P/P'\) and \(P'/P''\) are torsion-free, so there exists a complement of \(P''\) in \(P'\) and a complement of \(P'\) in \(P\) . Consequently, there exists a basis \((z_1, z_2, \ldots, z_l)\) of the \(Z\) -module \(P\) and rational integers \(j, m, n\) such that \(u = jz_1, v = mz_1 + nz_2, j > 0, n > 0\) . Putting \(X_i = e^{z_i}\) for \(1 \leqslant i \leqslant l\) , we then have \(1 - e^u = 1 - X_1^j, 1 - e^v = 1 - X_1^m X_2^n\) . Let \(K\) be an algebraic closure of the field of fractions of \(A\) , so that \(A[P]\) can be identified with a subring of the ring \(B = K[X_1, \ldots, X_l, X_1^{-1}, \ldots, X_l^{-1}]\) . For any \(j\) th root of unity \(z, 1 - zX_1\) is extremal in

\[
\mathrm{K} \left[ \mathrm{X} _ {1}, \dots , \mathrm{X} _ {l} \right];
\]

moreover, the ideal generated by \(1 - z\mathrm{X}_1\) contains no monomial in the \(\mathbf{X}_i\) . We conclude that the ideal \((1 - z\mathrm{X}_1)\mathrm{B}\) of \(\mathbf{B}\) is a prime ideal of height 1 (Commutative Algebra, Chap. VII, § 1, no. 6), hence that \(1 - zX_{1}\) is extremal in B. The extremal factors of \(1 - X_{1}^{j}\) in B are thus of the form \(1 - zX_{1}\) . Now none of these factors divide \(1 - X_{1}^{m}X_{2}^{n}\) in B (for the homomorphism \(f\) from B to B such that \(f(X_{1}) = z^{-1}\) , \(f(X_{i}) = X_{i}\) for \(i \geqslant 2\) , satisfies

\[
f (1 - z X _ {1}) = 0 \quad \text { and } \quad f (1 - X _ {1} ^ {m} X _ {2} ^ {n}) = 1 - z ^ {- m} X _ {2} ^ {n} \neq 0).
\]

Thus, \(1 - X_1^j\) and \(1 - X_1^m X_2^n\) are relatively prime in B. Consequently, any common divisor of \(1 - X_1^j\) and \(1 - X_1^m X_2^n\) in A{[}P{]} is invertible in B and so, up to multiplication by an element of the form \(X_1^{k_1}X_2^{k_2}\ldots X_l^{k_l}\) , is equal to an element \(a\) of A; in other words, \(a\) divides 1 in A, hence is invertible in A. Thus, finally, \(1 - X_1^j\) and \(1 - X_1^m X_2^n\) are relatively prime in A{[}P{]}.

